\chapter*{Abstract}

\section*{English}
This thesis proposes a learning-guided system for 3D nesting and transport of convex items that optimizes utilization under non-overlap, transport and container-bound constraints. It aims to explore use of neurogenesis in order to appraise its viability for this class of problems.

It proposes state exploration using heuristics to enumerate feasible loading plans combined with neural scoring models and architectures to rank candidate solutions. Model architectures are evolved via neurogenesis.

Items, containers, and candidates are represented with compact local and global features capturing size, capacity usage, and transport attributes. We survey the state of the art and build a reproducible benchmark. Evaluation reports utilization, bin count, cost, and runtime.

\textbf{Keywords:} nesting, neurogenesis, artificial intelligence, bin packing

\textbf{Supervisor:} Ing. Mgr. Ladislava Smítková Janků, Ph.D.

\section*{Abstrakt (Czech)}
Tato práce navrhuje systém pro 3D nesting a transport konvexních objektů řízený učením, který optimalizuje využití prostoru při dodržení vymezení jako je nepřekrývání, přepravních požadavků a atributy kontejneru. Cílem je prozkoumat použití neurogeneze a posoudit její vhodnost pro tuto třídu úloh.

Navrhuje prozkoumávání stavového prostoru pomocí heuristik, které enumerují proveditelné plány nakládky, v kombinaci s neuronovými skórovacími modely a architekturami pro řazení kandidátních řešení. Architektury modelů jsou vyvíjeny evolučně pomocí neurogeneze.

Položky, kontejnery a kandidáti jsou reprezentováni kompaktními lokálními i globálními vlastnostmi zachycujícími rozměry, využití kapacity a přepravní atributy. Posuzuje současné state-of-the-art řešení této třídy problému a navrhuje reprodukovatelný benchmark. Vyhodnocení bude uvádět míru využití, počet binů, náklady a dobu běhu.

\textbf{Klíčová slova:} polohování, neurogeneze, umělá inteligence, bin packing

\textbf{Překlad názvu:} Neurogeneze neuronových sítí pro 3D nesting
